\begin{textsample}{NEG Dim 3 – human – Score 1.00 – mg/mg\_ef\_ciencias\_humas\_intro\_1\_p3.txt}  \label{ex:f3_neg_014}
O ensino de Ciências Humanas, baseado na noção de tempo, espaço, movimento e interculturalidade, categorias básicas da área, deve estar associado à crítica sistemática à ação humana, às relações sociais e de poder e especialmente à produção de conhecimentos e saberes, frutos de diferentes circunstâncias históricas, espaços geográficos, conhecimentos religiosos e das filosofias de vida.
Enquanto Ciências Humanas, o ensino de Geografia, História e Ensino Religioso, deve estimular os estudantes a desenvolverem uma melhor compreensão do mundo, não só favorecendo o desenvolvimento autônomo de cada indivíduo, como também tornando -os aptos a uma intervenção mais responsável no espaço em que vivem.

% matched lemmas: 
\end{textsample}
